%==============================================================================
\part{The Setting}
%==============================================================================

\section{Introduction}
\label{sec:intro}

\subsection{The script that everybody writes}

Consider a question of the kind that arises daily in metabolic research: which
enzymes catalyse reactions consuming any member of a given chemical class, and
of those, which occur in pathways that also contain a second class. No single
public resource answers it. The reaction membership lives in one database, the
enzyme--reaction association in a second, the pathway membership in a third,
and the chemical class hierarchy in a fourth.

The universal solution is a script. In outline:

\begin{lstlisting}[caption={The script everybody writes. Not a strawman.},label=lst:script]
compounds = sparql(ENDPOINT_A, TEMPLATE_A % class_iri)
ids       = [row["c"] for row in compounds]
enzymes   = []
for chunk in batches(ids, 50):
    q = TEMPLATE_B % " ".join("<%s>" % i for i in chunk)
    enzymes += sparql(ENDPOINT_A, q)
kegg_ids  = [xref(e) for e in enzymes]
pathways  = [rest_get(ENDPOINT_C, "/link/pathway/%s" % k) for k in kegg_ids]
\end{lstlisting}

Every practitioner recognises \cref{lst:script}, and every practitioner
dislikes it. The dislike is usually expressed aesthetically. It is not an
aesthetic problem. \Cref{lst:script} exhibits, in eight lines, five distinct
and independently serious defects.

\begin{enumerate}[label=(D\arabic*),leftmargin=3em]
\item \textbf{The interpolation is textual.} Line~5 constructs a query by
string concatenation. Whether the resulting string means what the author
intended depends on the target language's parse of a string the author never
sees in full. \Cref{sec:interp} exhibits a documented case in which two
spellings of one intent, differing in a single line, returned $2$ and $397$
against the same service on the same day.
\item \textbf{Failure is one bit wide.} If \texttt{compounds} comes back empty,
the script cannot distinguish: the class has no members; the endpoint timed out
and returned an empty result rather than an error; the query used a construct
the endpoint does not support and the endpoint degraded silently; the class
identifier was misspelled. Each demands a different repair, and the script sees
one empty list.
\item \textbf{The failure of step $n$ is attributed to step $n$.} If
\texttt{pathways} is empty, the proximate cause is very often that
\texttt{kegg\_ids} was short, because \texttt{xref} is a partial function and
silently dropped two-fifths of its input. The script reports a pathway problem
and the analyst debugs the wrong stage.
\item \textbf{The budget is unstated and unallocated.} There is a rate limit on
\texttt{ENDPOINT\_C}, a timeout on \texttt{ENDPOINT\_A}, and no representation
of either anywhere in the program. The script discovers them by failing.
\item \textbf{The control flow is invisible to the data layer.} The loop over
\texttt{batches} exists because a query language has no way to say ``and
then''. The batching constant $50$ is a magic number encoding an unstated
belief about request-size limits at a service the script does not name.
\end{enumerate}

These are not five instances of sloppiness. They are five consequences of one
structural decision --- that control flow lives in a host language and data
access lives in a query language, and neither can see the other.

\subsection{Why the obvious repair fails}

The obvious repair is to make the query language expressive enough to hold the
whole computation: one large query with unions, optional patterns, subselects
and federated service clauses. This is what the SPARQL~1.1 federation extension
was designed for \citep{Prudhommeaux2013Federated,Harris2013}, and it fails for
two reasons that are not implementation defects.

First, it is frequently not available. Of the four sources in the motivating
question, only two speak SPARQL at all. A pathway resource exposing a flat-file
REST interface \citep{Kanehisa2000,Kanehisa2023} cannot be reached by a service
clause, because there is no SPARQL protocol endpoint to address. The federation
extension federates SPARQL endpoints, and the biological data landscape is not
made of SPARQL endpoints.

Second, where it is available, it concentrates the defects rather than removing
them. A single query containing every stage of the computation is a single
query whose failure is still one bit wide (D2), whose intermediate results are
still invisible (D3), and which is now large enough that the interpolation
defect (D1) becomes very hard to see --- as \cref{sec:interp} demonstrates on a
query of eight lines.

A third repair, the workflow system \citep{Wolstencroft2013,Koster2012,Berthold2009},
supplies exactly the control flow (D5) that the query language lacks. It leaves
(D1)--(D4) untouched, because a workflow node that issues a query issues a
string that its author assembled, and the workflow engine does not know what
the string means.

\subsection{What we build instead}

We separate the two things the script conflates. A \emph{plan} is a sequence of
\emph{steps}; a step names a source, an abstract pattern to send it, and a
binding to receive the result; the plan language supplies the control flow and
the source adapters supply the concrete requests. The user writes the plan.
Nobody writes a query.

\begin{principle}[Queries are leaves]
\label{prin:leaves}
No construct of the plan language denotes a graph pattern, a triple, a
relational clause, or a URL. Every such object is produced by \emph{lowering},
from an abstract step together with a source declaration, and is not
addressable by the plan author.
\end{principle}

\Cref{prin:leaves} is what makes the language portable across backends that
share no query model. Its cost is stated immediately: a plan cannot express a
request that its source adapters cannot generate, and \cref{sec:limits} records
what that excludes.

\begin{principle}[A step returns a verdict, not rows]
\label{prin:verdict}
The value of a step is a pair (verdict, payload). The payload is a result set
only when the verdict is $\vans$. Every other verdict carries a diagnosis and
no result set.
\end{principle}

\Cref{prin:verdict} is what makes the language diagnostic. Its content is
\cref{sec:verdicts}, and the reason it is not a mere convention is
\cref{cor:onebit}: an interface that returns rows, and signals failure by
returning zero of them, cannot represent the distinction at all, whatever
discipline is imposed on its callers.

\begin{principle}[Refusal is a first-class outcome]
\label{prin:refusal}
A plan whose declared requirements exceed what its sources can supply, or whose
budget falls below the minimum the plan needs, terminates without issuing any
request and reports what was missing.
\end{principle}

A system with no refusals declines nothing, and therefore claims everything it
is asked. \Cref{prin:refusal} is the operational form of that observation, and
\cref{thm:static} is the result that makes it enforceable statically rather
than by convention.

\subsection{Contributions and non-contributions}

We claim: a source model under which heterogeneous backends compose
(\cref{sec:model}); a capability calculus deciding compilability before
execution (\cref{thm:static}); a six-valued verdict algebra with a proof that
the conventional interface conflates its members (\cref{thm:six},
\cref{cor:onebit}); exact laws for retention under partial identifier
translation (\cref{thm:retention}) with a proof that reordering does not
recover loss (\cref{thm:reorder}) and that two translation routes may disagree
without either being wrong (\cref{thm:route-extent}); a budget allocation with
a single shadow price (\cref{thm:allocation}) and a proof that sufficient total
budget is not sufficient (\cref{prop:necessary-not-sufficient}); a proof that
structural interpolation eliminates a defect family that textual interpolation
admits (\cref{thm:interpolation}); and a characterisation of safe re-execution
(\cref{thm:idempotent}).

We do not claim: that plans written in this language return biologically
correct answers; that the capability declarations of \cref{sec:model} are
verified rather than asserted; that the retention laws of
\cref{sec:translation} have been measured against real cross-reference tables;
or that the prototype of \cref{sec:proto} has been exercised against any live
service. \Cref{sec:limits} enumerates these and others, and
\cref{rem:honesty-assumption} states the single assumption on which the largest
part of the development rests.

\subsection{Conventions}

$\Univ$ denotes a countable universe of identifiers. A \emph{namespace} is a
subset of $\Univ$; namespaces are pairwise disjoint, which is a modelling
decision and not a fact about the world (\cref{rem:disjoint}). Result sets are
finite. We write $f : X \rightharpoonup Y$ for a partial function and
$\operatorname{dom} f$ for its domain of definition. All complexity statements
count \emph{requests} unless stated otherwise, since request count is what a
public service rate-limits and what a plan can control; wall-clock time is not
modelled, and \cref{sec:limits} records the consequence.

\section{Related work and where this sits}
\label{sec:related}

Federated query processing over distributed data has a substantial literature
concerned principally with \emph{source selection} and \emph{join ordering}:
given a query and a set of endpoints, decide which endpoints can contribute and
in what order to contact them \citep{Schwarte2011,Gorlitz2011,Acosta2011,Saleem2016,Rakhmawati2013}.
That work assumes a single query language across sources and optimises within
it. Our setting is complementary: the sources do not share a language, the
decomposition is written by the user rather than inferred, and the object of
study is what a step reports when it fails rather than how fast it succeeds.
The two concerns are orthogonal, and a planner of the kind those papers
describe could in principle choose among the routes that \cref{sec:translation}
shows to be inequivalent --- but it would need a cost model over partial maps
that the literature does not supply.

Mediator and wrapper architectures \citep{Wiederhold1992,GarciaMolina1997,Haas1997}
established the pattern of a per-source adapter presenting a common interface,
and capability-based query planning \citep{Levy1996,Papakonstantinou1995,Florescu1999,Yerneni1999}
formalised sources that can answer only some queries, binding-pattern
restrictions being the classical instance. Our capability sets
(\cref{def:capability}) are in that lineage, and \cref{thm:static} is a
static-check analogue of their planning-time feasibility tests. The difference
lies in what happens on failure: those systems treat an uncoverable query as a
planning failure, and we treat it as a typed verdict that the plan can catch
and route around, which is the difference between a compiler error and an
exception.

Scientific workflow systems \citep{Wolstencroft2013,Berthold2009,Koster2012,DiTommaso2017}
provide the control flow we require, and are widely used for exactly the
motivating question. They are host languages with library calls, so the requests
they issue are strings assembled by the workflow author, and defects
(D1)--(D4) survive unchanged. A workflow system sequences steps; it does not
know what a step means, and so cannot report why one came back empty.

Identifier translation across biological namespaces has its own literature and
its own well-documented pathologies \citep{Chen2005,VanIersel2010,Wimalaratne2018,Juty2012,Ison2016}.
The standard framing is pairwise coverage: what fraction of namespace $A$ maps
into namespace $B$. \Cref{sec:translation} takes the \emph{composition} of
several such maps as the object, and derives what survives a chain --- a
quantity a plan actually depends on and which is not determined by the pairwise
coverages alone (\cref{prop:pairwise-insufficient}).

Finally, the observation that an empty result and a failure are different, and
that conflating them corrupts downstream reasoning, is old in database theory:
it is the null-value problem \citep{Codd1979,Imielinski1984,Libkin2016} in a new
setting. \Cref{sec:verdicts} is a six-valued refinement specific to federation,
and \cref{prop:starve-reachable} identifies a value with no analogue in the
single-database case.
